\documentclass[12pt]{article}
\usepackage[utf8]{inputenc}
\usepackage[T1]{fontenc}
\usepackage[margin=1in]{geometry}
\usepackage{setspace}
\usepackage{xcolor}
\usepackage[colorlinks=true, urlcolor=blue!60!black, linkcolor=blue!60!black]{hyperref}
\setlength{\parindent}{2em}
\setlength{\parskip}{1.25em}
\renewcommand{\baselinestretch}{1.0}
\begin{document}
\begin{center}
{\LARGE \textit{Boumediene v. Bush}, 553 U.S. 723 (2008)\par}
\vspace{0.8em}
{\large Case Brief\par}
\end{center}
\vspace{1em}

\section{Memory Jogger}
Holding 5--4 that the Suspension Clause has full effect at Guantanamo Bay, the Court held \S~7 of the Military Commissions Act unconstitutional as applied to these detainees because the Detainee Treatment Act's review of CSRT determinations was not an adequate and effective substitute for habeas corpus (\href{https://supreme.justia.com/cases/federal/us/553/723/}{Justia}).

\section{Detailed Case Facts}
Following the September 11, 2001 attacks, Congress enacted the Authorization for Use of Military Force (AUMF), empowering the President to use necessary and appropriate force against those who planned, authorized, committed, or aided the attacks. Petitioner Lakhdar Boumediene (an Algerian-born Bosnian citizen seized in Bosnia and Herzegovina) and other foreign nationals captured in Afghanistan or elsewhere abroad were transferred to the U.S. Naval Station at Guantanamo Bay, Cuba---territory over which Cuba retained formal, technical \textit{de jure} sovereignty under a 1903 lease agreement, while the United States exercised complete and indefinite \textit{de facto} jurisdiction and control.

Each petitioner denied being an enemy combatant. After \textit{Rasul v. Bush}, 542 U.S. 466 (2004), held that the federal habeas statute (28 U.S.C. \S~2241) extended to Guantanamo, the Defense Department established Combatant Status Review Tribunals (CSRTs) to determine whether detainees were ``enemy combatants.'' Under CSRT procedures, detainees had no right to legal counsel, were assigned only a non-lawyer ``Personal Representative,'' could not view classified evidence forming the basis of their detention, faced a rebuttable presumption in favor of the Government's evidence, and had no bar against hearsay.

Congress then enacted the Detainee Treatment Act of 2005 (DTA), attempting to strip federal habeas jurisdiction over Guantanamo detainees and channel limited review of CSRT determinations exclusively to the D.C. Circuit. After \textit{Hamdan v. Rumsfeld}, 548 U.S. 557 (2006), construed the DTA not to strip jurisdiction over pending habeas petitions, Congress enacted the Military Commissions Act of 2006 (MCA). Section 7 of the MCA amended 28 U.S.C. \S~2241(e) to strip federal courts of habeas jurisdiction over ``all cases, without exception,'' pending or future, brought by noncitizens detained as enemy combatants.

\section{Procedural History}
The D.C. Circuit (476 F.3d 981 (2007)) held that MCA \S~7 stripped jurisdiction over petitioners' pending cases and did not violate the Suspension Clause because noncitizens detained outside sovereign U.S. territory had no constitutional right to habeas corpus. After initially denying certiorari, the Supreme Court reconsidered and granted review on June 29, 2007.

\section{Judicial Votes}
The Supreme Court \textbf{reversed and remanded} the D.C. Circuit by a \textbf{5--4} vote. The case was argued on December 5, 2007, and decided on June 12, 2008 (\href{https://supreme.justia.com/cases/federal/us/553/723/}{Justia}; \href{https://www.law.cornell.edu/supct/html/06-1195.ZO.html}{Cornell LII}):

\begin{itemize}
    \item \textbf{Majority Opinion:} Justice Kennedy delivered the opinion of the Court, joined by Justices Stevens, Souter, Ginsburg, and Breyer.
    \item \textbf{Concurring Opinion:} Justice Souter filed a concurring opinion, joined by Justices Ginsburg and Breyer.
    \item \textbf{Dissenting Opinions:} Chief Justice Roberts filed a dissenting opinion, joined by Justices Scalia, Thomas, and Alito. Justice Scalia filed a dissenting opinion, joined by Chief Justice Roberts and Justices Thomas and Alito.
\end{itemize}

\section{Holding}
\textbf{Questions Presented:}

\begin{enumerate}
  \item Does MCA \S~7 strip federal courts of statutory habeas jurisdiction over pending petitions filed by Guantanamo Bay detainees?
  \item Do foreign nationals designated as enemy combatants and detained at Guantanamo Bay have the constitutional privilege of habeas corpus under the Suspension Clause?
  \item If the Suspension Clause applies, does the DTA's mechanism for D.C. Circuit review of CSRT determinations provide an adequate and effective substitute for habeas corpus, or does MCA \S~7 operate as an unconstitutional suspension of the writ?
\end{enumerate}

\textbf{Holding:} First, MCA \S~7 unambiguously strips federal courts of statutory habeas jurisdiction over pending Guantanamo petitions. Second, the Suspension Clause has full effect at Guantanamo Bay under a functional, non-formalistic extraterritoriality test, and petitioners are not barred from invoking it by their noncitizen status or enemy-combatant designation. Third, because DTA/CSRT review is not an adequate and effective substitute for habeas corpus---and Congress did not purport to invoke a formal Suspension Clause emergency---MCA \S~7 operates as an unconstitutional suspension of the writ.

\section{Analysis of Opinions}
\subsection{Majority Opinion (Justice Kennedy)}

Writing for a five-Justice majority, Justice Kennedy first agreed with the Government that MCA \S~7 clearly eliminated statutory habeas jurisdiction over pending cases, forcing the Court to confront the Suspension Clause directly. Tracing the history of the Great Writ from Magna Carta and \textit{Darnel's Case} (1627) to the Framing, Justice Kennedy emphasized that the Suspension Clause is not merely an individual right; it is an indispensable separation-of-powers mechanism designed to prevent executive detention without judicial oversight.

\textbf{1. Functional Extraterritorial Reach of the Suspension Clause:} Finding the 1789 common-law historical record inconclusive on territories like Guantanamo, the Court rejected the Government's bright-line rule that \textit{de jure} sovereignty marks the outer boundary of the Constitution. Synthesizing the \textit{Insular Cases}, \textit{Reid v. Covert}, 354 U.S. 1 (1957), and \textit{Johnson v. Eisentrager}, 339 U.S. 763 (1950), Justice Kennedy held that ``questions of extraterritoriality turn on objective factors and practical concerns, not formalism.'' Allowing the political branches to evade the Constitution simply by leasing foreign territory would let them ``switch the Constitution on or off at will.'' The Court articulated a three-factor functional test to determine the Suspension Clause's reach:
\begin{enumerate}
  \item \textbf{Citizenship, Status, and Adequacy of Process:} Unlike the \textit{Eisentrager} prisoners (who were convicted by military commissions after adversarial trials), the Guantanamo detainees contested their enemy-combatant status after far more limited, non-adversarial CSRT hearings.
  \item \textbf{Nature of the Sites of Apprehension and Detention:} Although petitioners were captured abroad, Guantanamo Bay had been under complete, exclusive, and indefinite U.S. control for over a century; Cuban law did not apply there.
  \item \textbf{Practical Obstacles:} Unlike post-World War II occupied Germany in \textit{Eisentrager}, Guantanamo is not an active theater of war, and adjudicating habeas petitions poses no significant military friction or mission disruption.
\end{enumerate}

\textbf{2. Inadequacy of the DTA/CSRT Substitute:} Even when Congress strips statutory habeas jurisdiction, it does not violate the Suspension Clause if it provides an ``adequate and effective substitute'' (\textit{Swain v. Pressley}, 430 U.S. 372 (1977)). Here, however, the combination of CSRT procedural defects (no counsel, limited access to classified allegations, unrestricted hearsay) and the DTA's narrow scope of D.C. Circuit review created an unacceptable risk of erroneous, potentially lifelong detention. Specifically, to constitute an adequate substitute where detention rests on executive determination rather than criminal trial, the reviewing court must have authority to (a) correct factual and legal errors, (b) assess the sufficiency of the Government's evidence, (c) admit and consider relevant exculpatory evidence discovered outside the CSRT record, and (d) order the detainee's conditional release. Because the DTA barred detainees from introducing newly discovered exculpatory evidence and left the court's release power uncertain, MCA \S~7 was unconstitutional. Finally, given that petitioners had already been detained for six years, the Court declined to require prudential exhaustion of DTA review before permitting district-court habeas proceedings.

\subsection{Concurring Opinion (Justice Souter)}

Justice Souter emphasized that \textit{Boumediene} was no sudden departure from precedent, because \textit{Rasul v. Bush} had already recognized the United States' plenary control over Guantanamo. He also stressed the sheer length of the detainees' six-year confinement without meaningful judicial review to rebut the Chief Justice's claim that the Court acted prematurely.

\subsection{Dissenting Opinion (Chief Justice Roberts)}

Chief Justice Roberts argued that the majority struck down ``the most generous set of procedural protections ever afforded aliens detained by this country as enemy combatants'' without even allowing the D.C. Circuit to test whether the DTA could be construed to afford adequate review. He criticized the majority for bypassing traditional constitutional avoidance and exhaustion principles, leaving detainees with years of undefined habeas litigation rather than prompt DTA adjudication.

\subsection{Dissenting Opinion (Justice Scalia)}

Writing in sharp originalist and separation-of-powers terms, Justice Scalia argued that \textit{Eisentrager} established a clear territorial bar against extending constitutional habeas rights to enemy aliens held outside sovereign U.S. territory, and that 1789 common law confirmed the writ never ran to non-subjects abroad. He denounced the majority's three-factor test as a manipulable, ``judicially brainstormed'' formula that usurped war-making authority from the political branches and warned that forcing military commanders to litigate battlefield detentions in civilian courts would have disastrous national-security consequences.

\subsection{Subsequent Developments}

\begin{itemize}

\item \textbf{Limits on Extraterritorial Reach:} Applying \textit{Boumediene}'s three-factor test, the D.C. Circuit held in \href{https://caselaw.findlaw.com/court/us-dc-circuit/1524965.html}{\textit{Al Maqaleh v. Gates}, 605 F.3d 84 (D.C. Cir. 2010)} (and reaffirmed in \href{https://law.justia.com/cases/federal/appellate-courts/cadc/12-5404/12-5404-2013-12-24.html}{\textit{Al Maqaleh v. Hagel}, 738 F.3d 312 (D.C. Cir. 2013)}), that the Suspension Clause did \textit{not} extend to noncitizen detainees held at Bagram Airfield in Afghanistan because Bagram lay inside an active theater of war where the United States lacked Guantanamo's century-long, exclusive control.

\item \textbf{Cabining the Suspension Clause in Immigration Law (\textit{Thuraissigiam}):} In \href{https://www.law.cornell.edu/supremecourt/text/19-161}{\textit{Department of Homeland Security v. Thuraissigiam}, 591 U.S. 103 (2020)} (\href{https://supreme.justia.com/cases/federal/us/591/19-161/}{Justia}), the Supreme Court (7--2 on the Suspension Clause) sharply limited \textit{Boumediene}'s expansionist potential. Upholding 8 U.S.C. \S~1252(e)(2)'s bar on habeas review of expedited-removal credible-fear determinations, Justice Alito held that the Suspension Clause protects only the writ ``as it existed in 1789''---a remedy to secure \textit{release from unlawful executive custody}, not a vehicle to obtain asylum or lawful admission into the United States.

\end{itemize}

\section{Examples: Future Applications}
\begin{enumerate}

\item \textbf{Same-Side (Suspension Clause Applies): A Leased Atoll Outside Any War Zone.} Suppose the United States signs a 99-year lease with a Pacific island nation granting ``exclusive U.S. administrative control'' over a remote atoll used to detain noncitizens suspected of cyber-terrorism, and Congress strips federal habeas jurisdiction while substituting an internal military review board that bars counsel and newly discovered exculpatory evidence. Under \textit{Boumediene}'s three factors---(1) contested status with weak process, (2) long-term exclusive \textit{de facto} U.S. control outside an active war zone, and (3) minimal practical obstacles---the Suspension Clause applies, and the jurisdiction-stripping statute is an unconstitutional suspension unless an actual rebellion or invasion requires suspension for public safety, or Congress provides an adequate judicial substitute.

\item \textbf{Same-Side (Inadequate Substitute): Judicial Review Without the Power to Order Release.} After \textit{Boumediene}, Congress again strips district courts of habeas jurisdiction over Guantanamo detainees and substitutes review in the D.C. Circuit. The new statute lets detainees have counsel, see the evidence against them, and present newly discovered exculpatory evidence, and it authorizes the court to declare a detention unlawful. But it provides that the court may only recommend release to the Secretary of Defense, who retains final authority. \textit{Boumediene} held that an adequate substitute for habeas must include, at a minimum, the power to order the conditional release of a person unlawfully detained. Because the reviewing court cannot order release, the substitute remains inadequate and the statute operates as an unconstitutional suspension of the writ.

\item \textbf{Opposite-Side (Suspension Clause Does Not Apply): A Forward Operating Base in an Active Combat Zone.} A foreign national captured and held by U.S. forces at a temporary forward operating base inside an active foreign combat zone would likely lack a constitutional habeas right if the United States lacks permanent exclusive control and judicial review would create severe practical and military obstacles.

\item \textbf{Opposite-Side (Outside the Core of the Writ): Expedited Removal at the Border.} A noncitizen apprehended 25 yards inside the U.S. border who seeks habeas review not to obtain release from custody, but to compel a second asylum hearing and remain in the country, cannot invoke \textit{Boumediene} because such relief falls outside the historical 1789 core of the writ.

\item \textbf{Fence-Sitter (Close Call): A Detention Facility on an Allied Nation's Military Base.} Under a long-term status-of-forces agreement, the United States operates a detention facility on a military base in an allied country far from any active hostilities, holding noncitizens captured in a distant conflict; the host nation retains sovereignty and its law formally applies on the base, though U.S. forces run the facility. Congress bars habeas review and provides only military status hearings. Two \textit{Boumediene} factors favor the detainees: the process for determining their status is weak, and there is no war-zone friction. But the site factor points the other way, because unlike Guantanamo, the United States lacks complete and indefinite control, and a court order might create friction with the host government. How courts weigh that mix is uncertain under \textit{Boumediene}'s functional test.

\end{enumerate}

\section{Critique}
\textbf{Indeterminacy of the Functional Test:} Justice Kennedy's three-factor test avoids formalistic evasion of the Constitution, but it provides lower courts little bright-line guidance on when U.S. military or diplomatic control over a foreign site crosses the threshold into \textit{de facto} sovereignty.

\textbf{Remedy Without Substantive Rules:} \textit{Boumediene} guaranteed Guantanamo detainees access to federal habeas courts but deliberately left open the substantive definition of ``enemy combatant,'' the burden of proof, and evidentiary rules---inviting years of D.C. Circuit litigation that ultimately narrowed detainee victories.

\textbf{Tension with \textit{Thuraissigiam}:} Compare \textit{Boumediene}'s living-functionalist approach to the Suspension Clause (rejecting 1789 territorial formalism as inconclusive) with \textit{Thuraissigiam}'s strict 1789 originalism---an acute methodological shift reflecting the Court's changed composition between 2008 and 2020.

\section{Key Quotations}
\begin{quote}
``To hold the political branches have the power to switch the Constitution on or off at will is quite another.''
\end{quote}
\noindent --- Justice Kennedy, for the Court, 553 U.S. at 765 (\href{https://supreme.justia.com/cases/federal/us/553/723/}{Justia}; \href{https://www.law.cornell.edu/supct/html/06-1195.ZO.html}{Cornell LII}).

\begin{quote}
``However, security subsists, too, in fidelity to freedom's first principles, chief among them being freedom from arbitrary and unlawful restraint and the personal liberty that is secured by adherence to the separation of powers.''
\end{quote}
\noindent --- Justice Kennedy, for the Court, 553 U.S. at 797 (\href{https://supreme.justia.com/cases/federal/us/553/723/}{Justia}; \href{https://www.law.cornell.edu/supct/html/06-1195.ZS.html}{Cornell LII}).

\begin{quote}
``Not the rule of law, unless by that is meant the rule of lawyers, who will now arguably have a greater role than military and intelligence officials in shaping policy for alien enemy combatants. And certainly not the American people, who today lose a bit more control over the conduct of this Nation's foreign policy to unelected, politically unaccountable judges.''
\end{quote}
\noindent --- Chief Justice Roberts, dissenting, 553 U.S. at 826 (\href{https://www.law.cornell.edu/supct/html/06-1195.ZD.html}{Cornell LII}; \href{https://supreme.justia.com/cases/federal/us/553/723/}{Justia}).

\begin{quote}
``Today, for the first time in our Nation's history, the Court confers a constitutional right to habeas corpus on alien enemies detained abroad by our military forces in the course of an ongoing war. \ldots{} The Nation will live to regret what the Court has done today.''
\end{quote}
\noindent --- Justice Scalia, dissenting, 553 U.S. at 826--27, 850 (\href{https://www.law.cornell.edu/supct/html/06-1195.ZD1.html}{Cornell LII}; \href{https://supreme.justia.com/cases/federal/us/553/723/}{Justia}).


\end{document}
